% Part: propositional-logic
% Chapter: propositional-logic 
% Section: soundness

\documentclass[../../../include/open-logic-section]{subfiles}

\begin{document}

\olfileid[pt]{pl}{prp}{snd}

\olsection{Correção da lógica proposicional}

\begin{thm}[Correção] 
\ollabel{thm:soundness} 
Se $\Gamma\Proves!A$, então
$\Gamma \Entails !A$. 
\end{thm}

\begin{proof} 
Por indução sobre os teoremas. Se $!A$ é um axioma, então $\Entails
!A$ e, portanto, a fortiori, $\Gamma \Entails!A$. Do mesmo modo, se $!A \in \Gamma$,
$\Gamma \Entails!A$. Para o passo indutivo, suponha que $!A$ seja obtida por
\emph{modus ponens} de $!C$ e $!C\lif!A$; suponha também, como hipótese de
indução, que o teorema valha para $!C\lif !A$ e $!C$. O terceiro item
de \olref[sem]{prop:semanticalfacts} fornece o resultado desejado. 
\end{proof}

\begin{cor} 
Se $\Gamma$ é satisfatível, então $\Gamma$ é consistente. Logo,
a lógica proposicional é consistente. 
\end{cor}

\end{document}
